
%-------------------------------------------------------------
% Kushagra Petwal - Resume (ATS-optimized, Google XYZ bullets)
% Compile with: pdflatex Kushagra_Petwal_Resume.tex
%-------------------------------------------------------------
\documentclass[letterpaper,10pt]{article}

\usepackage[T1]{fontenc}
\usepackage[scaled=0.92]{helvet}
\renewcommand{\familydefault}{\sfdefault}

\usepackage[letterpaper,top=0.4in,bottom=0.35in,left=0.55in,right=0.55in]{geometry}
\usepackage{microtype}
\usepackage{xcolor}
\usepackage{titlesec}
\usepackage{enumitem}
\usepackage[hidelinks,
  pdftitle={Kushagra Petwal - Resume},
  pdfauthor={Kushagra Petwal},
  pdfsubject={Embedded Systems and Firmware Engineer Resume}]{hyperref}
\input{glyphtounicode}
\pdfgentounicode=1

\definecolor{AccentColor}{RGB}{0,32,96}

\pagestyle{empty}
\setlength{\parindent}{0pt}
\setlength{\parskip}{0pt}
\setlength{\emergencystretch}{2em}
\raggedright
\urlstyle{same}

% ---------- Section headings ----------
\titleformat{\section}[block]
  {\color{AccentColor}\bfseries\normalsize\raggedright}
  {}{0em}
  {\MakeUppercase}
  [{\vspace{0.5pt}\color{AccentColor}\titlerule[0.6pt]}]
\titlespacing*{\section}{0pt}{7pt}{1pt}

% ---------- Bullet lists ----------
\setlist[itemize]{
    leftmargin=1.1em,
    labelsep=0.5em,
    itemsep=1pt,
    parsep=0pt,
    topsep=1.5pt,
    partopsep=0pt,
    label=\textbullet
}

% ---------- Entry helpers ----------
% \entryhead{bold left}{right}{italic left}{italic right}
\newcommand{\entryhead}[4]{%
  \addvspace{3pt}%
  \noindent\textbf{#1}\hfill #2\par\nobreak
  \noindent{\small\textit{#3}}\hfill{\small\textit{#4}}\par\nobreak
  \vspace{1pt}}

% \projhead{project name}{tech stack}{right}
\newcommand{\projhead}[3]{%
  \addvspace{3pt}%
  \noindent\textbf{#1} ~|~ {\small\textit{#2}}\hfill #3\par\nobreak
  \vspace{1pt}}

% \entryline{bold}{italic}{right}
\newcommand{\entryline}[3]{%
  \addvspace{3pt}%
  \noindent\textbf{#1}, \textit{#2}\hfill #3\par\nobreak
  \vspace{1pt}}

% \skillline{category}{keywords}
\newcommand{\skillline}[2]{%
  \par\hangindent=1.1em\noindent\textbf{#1:} #2\par\vspace{1.2pt}}

\begin{document}

%---------- HEADING ----------
\begin{center}
  {\huge\bfseries\color{AccentColor} Kushagra Petwal}\\[3pt]
  {\small Pantnagar, Uttarakhand, India ~|~
  \href{mailto:kushpetwal13@gmail.com}{kushpetwal13@gmail.com} ~|~
  \href{tel:+918650110088}{+91 8650110088}}\\[1pt]
  {\small
  \href{https://www.linkedin.com/in/kushagra-petwal-66b537224/}{linkedin.com/in/kushagra-petwal-66b537224} ~|~
  \href{https://github.com/KAVATER}{github.com/KAVATER}}
\end{center}
\vspace{-10pt}

%---------- EDUCATION ----------
\section{Education}
\entryhead
  {G. B. Pant University of Agriculture and Technology}{2027}
  {Bachelor of Technology (B.Tech) in Electronics and Communication Engineering}
  {Pantnagar, Uttarakhand, India}

\entryhead
  {St. Thomas' College}{}
  {Class XII: 80\% \quad | \quad Class X: 89\%}
  {Dehradun, Uttarakhand, India}

%---------- WORK EXPERIENCE ----------
\section{Work Experience}
\entryhead
  {Embedded Software Intern}{Jul 2026 -- Sep 2026}
  {Sunfox, Dehradun, Uttarakhand, India}{\href{https://github.com/KAVATER/SUNFOX-2026-Projects}{github.com/KAVATER/SUNFOX-2026-Projects}}
\begin{itemize}
  \item Increased ECG sampling throughput 6x (100 to 600 samples per second, SPS) on an nRF52840-based ECG device by developing Zephyr RTOS firmware, enabling Bluetooth Low Energy (BLE) 2M PHY, and optimizing the codebase.
  \item Designed a Finite Impulse Response (FIR) filter in MATLAB and deployed it with ARM CMSIS-DSP, using FFT-based processing of ECG signal data and filter coefficients.
  % XYZ-TODO: add result, e.g. noise reduction in dB, filter order/cutoff, or cycles per sample
  \item Delivered low-flicker, real-time ECG waveform rendering on a 2.4-inch TFT display using 8-bit fonts and a sliding-window graph that maps display pixels to a circular buffer.
  % XYZ-TODO: add frame rate / refresh rate
  \item Configured the STM32WB55 as a BLE Central (scanning, connection establishment, command-triggered GATT notifications) via the STM32 Wireless SDK, and as a BLE Peripheral with multi-device connectivity for ECG data transmission.
  % XYZ-TODO: add number of simultaneous connections
\end{itemize}

\entryhead
  {Embedded Intern}{Jul 2025 -- Sep 2025}
  {Sunfox, Dehradun, Uttarakhand, India}{\href{https://github.com/KAVATER/Sunfox-Projects/tree/main/HeartRate\%20and\%20Spo2-2/Core/Src}{github.com/KAVATER/Sunfox-Projects}}
\begin{itemize}
  \item Developed AFE4490 analog front-end breakout-board drivers and algorithms for SpO2 and heart-rate monitoring.
  \item Removed baseline drift from the SpO2 signal by applying FIR high-pass filters.
  \item Built a custom USB protocol analyzer as a cost-effective alternative to commercial tools.
  % XYZ-TODO: add cost saved (INR) or supported protocols/sample rate
  \item Implemented a Kalman filter for the MPU6050 inertial measurement unit (IMU) and used the Discovery Board accelerometer to analyze the effect of motion on SpO2 signals.
\end{itemize}

%---------- PROJECTS ----------
\section{Projects}
\projhead
  {Drone Firmware}{ESP32, Arduino, I2C, PID, Kalman Filter}{\href{https://github.com/lab685/ESP32-Drone}{github.com/lab685/ESP32-Drone} ~|~ 2025}
\begin{itemize}
  \item Built an ESP32-based drone flight controller on the Arduino framework with a real-time PID loop and MPU6050 IMU over I2C.
  \item Improved motion estimation by implementing a Kalman filter on IMU data.
  % XYZ-TODO: add control-loop rate (Hz) or stability result
\end{itemize}

\projhead
  {Custom Embedded Linux for PocketBeagle 2}{Buildroot, OP-TEE, ARM}{2026}
\begin{itemize}
  \item Built a custom embedded Linux distribution for the PocketBeagle 2 using Buildroot, configuring the bootloader, Linux kernel, root filesystem, and cross-compilation toolchain for an ARM-based platform.
  \item Integrated OP-TEE (Open Portable Trusted Execution Environment) and configured the device tree and peripheral drivers to enable secure embedded execution.
  % XYZ-TODO: add boot time / image size / repo link
\end{itemize}

\projhead
  {Battle Bot}{Fusion 360, PrusaSlicer, Pico SDK}{2026}
\begin{itemize}
  \item Designed a 3D-printed battle bot in Fusion 360 and fabricated its components on a Prusa MK4 3D printer; placed 3rd in the Beetleweight category at Robo War, IIT BHU.
  \item Developed Pico SDK firmware that decodes FlySky transmitter inputs and implements watchdog-based fail-safe mechanisms.
\end{itemize}

%---------- TECHNICAL SKILLS ----------
\section{Technical Skills}
\vspace{1pt}
\skillline{Languages}{C (Embedded C), Java, MATLAB}
\skillline{Platforms}{STM32 , ESP32, Nordic nRF52840,NXP, PocketBeagle 2, ARM Cortex-M/A}
\skillline{RTOS and Embedded Linux}{Zephyr RTOS, FreeRTOS, Buildroot}
\skillline{Protocols and Wireless}{Bluetooth Low Energy (BLE), GATT, UART, SPI, I2C, CAN, USB, MQTT}
\skillline{Signal Processing and Control}{FIR Filters, FFT, ARM CMSIS-DSP, Kalman Filter, PID Control}
\skillline{GUI and Tools}{TouchGFX, LVGL, STM32 Wireless SDK, Pico SDK, GNU Make, Git, Fusion 360, PrusaSlicer}

%---------- LEADERSHIP ----------
\section{Leadership}
\entryline{ChairPerson}{The Robotics Club, Pantnagar}{}% XYZ-TODO: add years held
\begin{itemize}
  \item Organized a bot-building workshop for 100+ first-year students; mentored and supervised bot-building competitions.
  % XYZ-TODO: add number of teams / events
\end{itemize}

\entryline{Member}{Vimarsh: The Oratory Club, Pantnagar}{}% XYZ-TODO: add years
\begin{itemize}
  \item Organized and participated in multiple university-level debate competitions.
\end{itemize}

% %---------- AWARDS ----------
% \section{Awards}
% \vspace{2.5pt}
% \begin{itemize}
%   \item \textbf{3rd Place (Second Runner-Up), Beetleweight Battle Bot Competition}, Robo War, IIT BHU \hfill 2026
%   \item \textbf{Runner-Up (2nd Place), Maze Solving Competition}, The Robotics Club, Pantnagar \hfill 2024
% \end{itemize}

\end{document}
